\section*{测度的射影极限}
\phantomsection
\addcontentsline{toc}{section}{测度的射影极限}

这是一项特别着眼于概率论的需要而发展起来的理论。关于随机变量的有限序列 \(X_{1}, \ldots, X_{n}\) 的问题，原则上当该序列的律 \(P_{X}\) 已知时便告解决：它是 \(\mathbf{R}^{n}\) 上质量为 1 的一个正测度，使得同时得到诸不等式 \(a_{1} \leq X_{1} \leq b_{1}, \ldots, a_{n} \leq X_{n} \leq b_{n}\) 的概率等于 \(P_{X}(C)\)，其中 \(C\) 是闭长方体 \([a_{1}, b_{1}] \times \ldots \times [a_{n}, b_{n}]\)（位于 \(\mathbf{R}^{n}\) 中）。

实际上，测度 \(P_{X}\) 具有离散支集，或者具有关于勒贝格测度的密度。当涉及随机变量的无穷序列 \((X_{n})_{n\geq1}\) 时，部分序列的律 \(P_{n}\)（即 \((X_{1},\ldots,X_{n})\) 的分布律）一般对每个整数 \(n\geq1\) 都是已知的；这些给定的律满足一个相容性条件，它表达了如下事实：序列 \((P_{n})_{n\geq1}\) 是一个测度的射影子系。直到 1920 年前后，人们还以多少是隐含的方式，通过从有限情形的概率出发作“自然”的过渡到极限，来定义与无穷序列相联系的事件的概率；例如，人们会假定一局博弈终止的概率，是它至多在 n 步内终止的概率当 n 趋于无穷时的极限。自然，这样的理论很难说是自洽的，而且没有什么能排除“悖论”的出现：同一个概率会因用两种各自同样“自然”的程序中的这一种或那一种来计算而得到两个不同的估计。

施坦豪斯{[}293{]}似乎是第一个感到有必要（就掷硬币的情形）不仅考虑射影子系 \((P_{n})_{n\geq1}\) 而且考虑其极限的人。稍早一些，1919 年，丹尼尔{[}76 d{]}已经在一般情形证明了这类射影极限的存在性，\(^{7}\)但这一结果在欧洲似乎一直无人知晓。1933 年，柯尔莫哥洛夫在著作{[}183{]}中重新得到它，他在此书中提出了概率论的公理化概念。丹尼尔与柯尔莫哥洛夫的证明都使用了一个依赖于迪尼定理的紧性论证。

丹尼尔–柯尔莫哥洛夫定理在随机序列 \((X_{n})_{n\geq1}\) 的情形已无可挑剔，但从 1935 年起由柯尔莫哥洛夫、费勒与杜布所开展的对随机函数的研究却隐藏着完全另一量级的困难。例如考虑 R 的一个区间 T，它代表“随机过程”的观测时刻所成的集；可能轨道的全体是乘积空间 \(R^{T}\)，它被看作部分乘积 \(R^{H}\) 的射影极限，其中 H 跑遍 T 的有限子集所成的集；一般就得到一个测度的射影子系 \((\mu_{H})\)。柯尔莫哥洛夫定理确实给出了 \(R^{T}\) 上的一个测度，但它只定义在一个比博雷尔族小得多的族上。\(^{8}\) 柯尔莫哥洛夫构造的一个给出拓扑空间上测度的变体属于角谷{[}177{]}，此后曾被多次重新发现：把 \(\mu_{H}\) 看作 \(\overline{R}^{H}\) 上一个由 \(R^{H}\) 承载的测度，\(^{9}\)紧空间 \(E=\overline{R}^{T}\) 是诸有限乘积 \(\overline{R}^{H}\) 的射影极限，而 E 上的测度 \(\mu\) 可以定义为诸 \(\mu_{H}\) 的射影极限。但这一程序有一个严重的弊病；\(\overline{R}^{T}\) 的元素不具有任何正则性性质，使人能继续该过程的概率研究——甚至不能简单地消去寄生的值 \(\pm\infty\)（由 R 的紧化 \(\overline{R}\) 所引入）。补救的办法是把测度 \(\mu\)（\(\overline{R}^{T}\) 上的）限制到一个特定的子空间（例如布朗运动情形中的 \(C(T)\)）；根本的困难来自这样一个事实：一个函数空间，即使是通常类型的，也未必是 \(\mu\)-可测的（在 \(\overline{R}^{T}\) 中），甚至函数空间的选择本身也可能成问题。\(^{10}\)

1956 年，普罗霍罗夫迈出了决定性的一步，他的著作{[}254{]}对随机过程理论产生了决定性的影响。他把上面分析的文章中维纳所用的方法放进一个方便的公理化形式，从而建立了函数空间上测度的射影极限的一般存在性定理。

1947 年，博赫纳{[}24 b{]}引入了一类较窄的射影子系，它们由有限维实向量空间与满射线性变换组成。这种子系的射影极限以自然的方式等同于代数对偶 \(E^*\)——这里 \(E\) 是装备着弱拓扑 \(\sigma(E^*, E)\) 的一个实向量空间；相应的测度射影子系有一个极限，即测度 \(\mu\)，它只对一个比 \(E^*\) 的博雷尔族小得多的族有定义。博赫纳借其傅里叶变换（它是 \(E\) 上的一个函数）完全刻画了这类“前测度”。但在 \(E\) 上没有拓扑时这一结果几乎无法使用，这时必须考察把 \(\mu\) 看作拓扑对偶 \(E'\)（即 \(E\) 的拓扑对偶）上的测度的可能性。与此独立地，R. 福尔泰与 E. 穆里耶在设法把概率论的某些经典结果（大数定律、中心极限定理）推广到取值于一个巴拿赫空间的随机变量时，也凸现了傅里叶变换在这些问题中所起的基本作用。但实质性的进展直到 1956 年才取得，当时盖尔范德{[}125 c{]}提出，傅里叶变换的自然框架既不是巴拿赫空间也不是希尔伯特空间，而是核型弗雷歇空间。他猜想，这样一个空间上的每个正型连续函数都是其对偶上一个测度的傅里叶变换，这一结果不久便由明洛斯{[}222{]}建立。其重要性首先在于它适用于广义函数空间，而且几乎全部函数空间都是广义函数空间的博雷尔子空间（因而它构成一个比 \(\mathbf{R}^T\) 好得多的框架）。\(^{11}\) 随机广义函数理论是一个正全面展开的领域，我们只须请读者参阅盖尔范德与维连金的著作{[}126{]}即可。

我们刚才提到的关于射影极限的结果都用到了底空间上拓扑的存在性。可以问，在“抽象”测度的情形是否存在一个类似的理论。冯·诺伊曼早在 1935 年就证明了乘积测度在所有情形的存在性，但延森与安德森{[}290{]}发现的一个反例摧毁了每个测度射影子系都承认极限这一希望。人们找到了两种权宜之计：1949 年，C. 扬内斯库-图尔恰在存在适当分解的条件下建立了可数射影极限的存在性，\(^{12}\) 这是研究马尔可夫过程的一个非常有趣的结果；此外，人们已经看清，空间的拓扑只是通过紧子集的全体才介入。因此，很自然地要借助一个集合的子集的紧类这一概念，在一种抽象理论内部把这种状况公理化。这项工作由马尔切夫斯基（他借此建立了一个关于抽象射影极限的定理）与雷尔-纳德泽夫斯基（他研究测度的分解）于 1953 年完成。\(^{13}\)

